\subsection{Commercial dynamic panels: instruments, horizons and the effective calendar}
\label{exec19:dynamic-scope}

A funding team may observe hundreds of branches but only a few policy dates.
Its balance equation can contain lagged balances, branch-specific conditions,
quoted rates, common liquidity shocks and delayed measurements. The software
comparison must separate four questions: forecasting balances, estimating a
dynamic panel coefficient under instrument restrictions, tracing a supplied
structural shock, and evaluating a policy that changes the pricing rule.
These questions share data structures but do not share an automatic causal
interpretation. Commercial products already provide substantial estimation,
diagnostic and reporting capabilities. The proposed common interface needs
to preserve their scientific premises before it can compare or compose them.

The additional inspected originals are Stata's \texttt{xtabond},
\texttt{xtdpdsys} and \texttt{lpirf} manuals and EViews' panel-estimation
chapter. The Stata PDFs carry July 2026 creation metadata; the EViews page
displays a March 2025 update. All were acquired on 19 October 2026. These
dates identify the inspected material but do not authenticate the exact bytes
available at the 7 October cutoff or every feature in a licensed edition.
The comparison is therefore documentary and version-limited. It does not
report a paid-engine trial or a complete commercial-market census.
\citep{exec19_dynamic_stata_ab,exec19_dynamic_stata_system,
exec19_dynamic_stata_lp,exec19_dynamic_eviews_panel}

\paragraph{Difference and system estimation answer conditional moment questions.}
Write a scalar dynamic panel model as
\begin{equation}
 B_{it}=\rho B_{i,t-1}+\beta A_{it}+\nu_i+\epsilon_{it},
 \label{exec19:dynamic-model}
\end{equation}
where $B_{it}$ is a balance measure, $A_{it}$ a rate or action variable,
$\nu_i$ a branch component and $\epsilon_{it}$ an idiosyncratic disturbance.
Differencing removes $\nu_i$, but does not make the differenced lagged
outcome exogenous. Lagged levels can supply moment conditions only under
appropriate timing and error restrictions. A rate chosen in response to
current liquidity news must not be classified as strictly exogenous merely
because its database column precedes the outcome column. Predetermined,
endogenous and strictly exogenous declarations change the admissible
instrument lags and the resulting estimator.

The inspected \texttt{xtabond} manual explicitly targets many panels and
few periods and requires absence of autocorrelation in the idiosyncratic
errors. It distinguishes one-step estimation, two-step estimation and their
variance estimators. Its robust two-step option uses a Windmeijer correction;
that correction addresses the declared variance calculation, not the truth
of the instrument moments. The manual also warns that instrument counts grow
quadratically with the number of dates and that too many overidentifying
restrictions can perform poorly in small samples. It supplies lag limits
and serial-correlation diagnostics. A diagnostic nonrejection is evidence
about a test under its assumptions, not proof that a bank's assignment
mechanism satisfies every required restriction.~\citep{exec19_dynamic_stata_ab}

System estimation adds equations in levels with lagged differences as
instruments. The \texttt{xtdpdsys} manual motivates these moments when
lagged-level instruments become weak under persistence, but states an
additional initial-condition orthogonality restriction. Extra equations do
not supply extra valid information unless that restriction is defensible.
The source retains the many-panel/few-period setting and idiosyncratic-error
conditions. A common service should consequently record which moments were
used and why they are valid, rather than translate a request for a more
persistent model directly into a stronger estimator label.
\citep{exec19_dynamic_stata_system}

An exact authored example isolates the initial-condition issue. Set the
action term to zero in \eqref{exec19:dynamic-model}, take $\rho=1/2$, and
let $\nu_i,\eta_i,\epsilon_{i2},\epsilon_{i3}$ be independent fair
signs. Under initialization $B_{i1}=\nu_i$,
\begin{equation}
 \Delta B_{i2}=\tfrac12\nu_i+\epsilon_{i2},\qquad
 E[\Delta B_{i2}(\nu_i+\epsilon_{i3})]=\tfrac12.
 \label{exec19:dynamic-invalid}
\end{equation}
The proposed level-equation moment at date three is nonzero. By contrast,
if $B_{i1}=\nu_i/(1-\rho)+\eta_i=2\nu_i+\eta_i$,
\begin{equation}
 \Delta B_{i2}=-\tfrac12\eta_i+\epsilon_{i2},\qquad
 E[\Delta B_{i2}(\nu_i+\epsilon_{i3})]=0.
 \label{exec19:dynamic-valid}
\end{equation}
Both identities follow by expanding the products and using the independence
and zero means of the signs. The reference enumerates all sixteen states
for each initialization. This is an assumption check, not a reproduction of
a commercial estimator's error or a proof that the second initialization
qualifies every system-GMM moment. Increasing the number of branches cannot
turn the nonzero population moment in \eqref{exec19:dynamic-invalid} into
zero. A covariance correction cannot do so either.
\citep{exec19_dynamic_reference}

\paragraph{Instrument construction is part of the output contract.}
EViews' inspected chapter exposes first differences and orthogonal deviations
for removing individual effects. Its \texttt{@DYN} syntax expands
period-specific lag instruments; \texttt{@LEV} preserves an instrument in
levels within a transformed equation. These choices are substantively
different from leaving a regressor unchanged or treating a lag as universally
valid. For one endogenous series with unrestricted period-specific levels
from lag two backward, a balanced history of $T$ dates produces
\begin{equation}
 K(T)=\sum_{t=3}^{T}(t-2)=\frac{(T-1)(T-2)}{2}.
 \label{exec19:dynamic-instruments}
\end{equation}
At $T=20$ this is $171$, matching the illustrative count in the inspected
EViews chapter. The formula describes this particular uncollapsed instrument
construction, not every GMM configuration or rank after missingness. The
adapter should return the actual columns, lags, count and rank, along with
the equation transformation and one-step or two-step weighting choices.
\citep{exec19_dynamic_eviews_panel,exec19_dynamic_reference}

The same EViews page distinguishes covariance methods labeled for period
clustering, individual clustering and two-way clustering. It also states
that some GMM weights cease to be available after differencing or orthogonal
deviation. One displayed weight label appears alongside a description of
two-step or multistep use that requires edition-specific clarification; the
proposal should not resolve that ambiguity by silently choosing its own
interpretation. Nor should it infer that every covariance option described
for a panel equation is available for every transformed GMM model. The
selected executable specification must bind the actual permitted combination.
\citep{exec19_dynamic_eviews_panel}

\paragraph{A branch count is not a count of independent policy dates.}
Consider an authored descriptive law for a branch-level funding response,
\begin{equation}
 Z_{it}=\theta+U_t+E_{it},
 \label{exec19:dynamic-calendar}
\end{equation}
with independent mean-zero, unit-variance common-date shocks $U_t$ and
idiosyncratic shocks $E_{it}$. Distinct dates are independent in this example.
There are $N$ branches and $T$ dates, with all rows observed. The overall mean
satisfies
\begin{equation}
 \operatorname{var}(\bar Z)=\frac1T+\frac1{NT},
 \qquad
 V_{\mathrm{row}}=\frac2{NT},
 \label{exec19:dynamic-variance}
\end{equation}
where $V_{\mathrm{row}}$ is the value obtained by incorrectly declaring
all $NT$ rows independent. The first identity follows by writing the mean as
$\theta+T^{-1}\sum_t U_t+(NT)^{-1}\sum_{i,t}E_{it}$. The common
shock averages over dates and does not disappear as more branches are added.
At $N=3,T=2$, enumeration of all $2^{2+6}=256$ fair-sign histories gives
variance $2/3$, while the false row-independent calculation gives $1/3$.
\citep{exec19_dynamic_reference}

For $N=100,T=25$, the variances are $101/2500$ and $1/1250$, a ratio of
$101/2=50.5$. If one response unit is stipulated to mean USD~10,000 per
branch-period, the corresponding standard deviations of the mean are
approximately USD~2,010 and USD~283. These are exact population-variance
consequences of the assumed law, with square roots rounded for presentation.
They are not confidence intervals, estimated bank losses or a report that
any named package uses the incorrect calculation. A software option for
clustering by date is useful, but twenty-five dates do not automatically
justify its finite-sample calibration. Persistent shocks, unequal loadings,
missing dates and selection require further analysis beyond this example.
\citep{exec19_dynamic_reference}

For a bank-wide rate policy that changes only by calendar date, adding
saturated date effects can also remove the very rate variation used to
identify its coefficient. Time controls and common-shock handling must be
designed with the target, not appended as a generic repair. Branch-specific
exposure variation, a justified external instrument, a different policy
design or a longer independent calendar may supply useful information, but
each changes the assumptions or the estimand. The proposal cannot promise
that a larger table alone closes the information gap.

\paragraph{Local projections and structural shocks remain distinct.}
Stata's \texttt{lpirf} manual describes joint multistep regressions, dynamic
multipliers for declared exogenous variables, horizon and lag controls, and
OLS, robust and HAC covariance choices. Its \texttt{dfk} option changes
the residual covariance divisor and \texttt{small} changes the reference
statistics. Long horizons consume trailing observations and can greatly
reduce the usable sample. These are material comparison dimensions. They
do not establish that a pricing innovation is independent of liquidity news
or that a pointwise interval is a simultaneous guarantee over every reported
horizon.~\citep{exec19_dynamic_stata_lp}

The previously inspected Stata structural-VAR manual requires supplied
short-run or long-run restrictions and explicitly warns that convergence
can occur without structural identification. Its external-instrument VAR
manual distinguishes a single target shock from multiple shocks and states
relevance, exclusion and normalization requirements. Those originals remain
the appropriate evidence for their bounded claims. A local projection, a
restricted VAR and an instrumented shock response should be compared at the
same shock units, information clock and horizons before their numerical
responses are compared.~\citep{exec13_ct_stata_svar,exec13_ct_stata_ivsvar}

Other commercial families provide useful alternatives with different current
evidence. The inspected SAS VARMAX text describes Granger tests and a
lower-triangular dynamic representation; it does not establish the full
dynamic-panel procedure or every equation omitted by HTML text extraction.
IBM's inspected TCM overview explicitly defines its use of causal in
predictive Granger terms. Mplus' selected DSEM/RDSEM addendum describes
supplied longitudinal models, random effects and missing-covariate limits.
None should be relabeled as an equivalent dynamic-panel policy estimator
without a more specific contract.~\citep{exec13_ct_sas_causal,
exec13_ct_sas_structural,exec13_ct_ibm_tcm,exec13_ct_mplus81}

\paragraph{What a public comparison can establish next.}
The current product proposal should offer typed requests for a chosen
scientific target, instrument design, sampling calendar and uncertainty
statement. Commercial engines remain possible licensed comparators; this
review supplies no entitlement, current price or right to embed them.
MathWorks' selected official IRF route was inaccessible, earlier OxMetrics
manual routes remain unresolved, and a hosted open-source frontend still
needs a concrete vendor, edition and numerical backend. These are specific
evidence omissions, not findings that the products lack relevant features.

A public or generated benchmark can first freeze a valid short-panel law and
an intentionally invalid initial-condition law, a common-date law, and
weak-instrument or rank-deficient cases. Each engine must receive the same
rows, lags, transformation, normalization and retained horizons. Record
instrument counts, effective dates, convergence, diagnostic limitations and
all failed attempts. A refusal should retain its reason rather than become
a zero effect. Compare point estimates only when their targets agree and
compare uncertainty only under a stated joint or pointwise criterion. This
work can justify whether a later licensed prototype and bank experiment is
worth funding while leaving universal dependence, hidden-confounding and
institutional acceptance claims open.

\section{A purchasing recommendation must retain the scientific comparison}

The detailed studies give the acquisition argument a concrete basis. Existing
systems provide graph exploration, supplied-model fitting, latent-class
analysis, trial-specific effects and temporal estimators that the proposed
product should compare seriously. Their outputs answer different questions.
The observationally equivalent lending models, imperfect proxy, compliance
example, predictive lag and alternative structural orderings show why fit or
interface convenience cannot supply the missing identification information.
The panel examples add instrument validity, initialization and the effective
calendar to that same requirement.

The cost and capacity derivations apply only after those scientific objects
match. Model residency, online information, deadlines and shared validation
dependencies can change which licensed configuration is useful. Hypothetical
cost crossings and compatibility signals identify measurements that could
change a purchase; they do not establish actual savings. The defensible
recommendation is conditional reuse with a preserving interface and a matched
comparison of the complete workflow. The next chapter develops the
identification and counterfactual distinctions that such an interface must
retain.
